% Gerado automaticamente por src/reporting/tables.py — não editar à mão.
\begin{table}[htbp]
\centering
\caption{Desempenho do mesmo modelo final em duas partições de teste. A partição do TCC I foi reconstruída pelo corte 80/20 por índice de linha; a partição correta segue o protocolo do Capítulo de metodologia. Como o modelo, os hiperparâmetros e o treino são idênticos, a diferença é atribuível à composição do conjunto.}
\label{tab:efeito-composicao}
\small
\begin{tabular}{lrrlrrrr}
\toprule
Conjunto de teste & Partidas & Acurácia & IC 95\% & F1 macro & F1 Derrota & F1 Empate & F1 Vitória \\
\midrule
Teste correto (WP1) & 340 & 46.18 & [41.76; 50.59] & 43.44 & 51.79 & 22.69 & 55.83 \\
Teste do TCC I (reconstruído) & 527 & 54.55 & [51.04; 58.16] & 50.36 & 60.78 & 27.11 & 63.18 \\
\bottomrule
\end{tabular}
\end{table}
